\documentclass[mestrado,12pt,onehalfspacing]{UnB-PPGT}

% Metadados obrigatórios para teste
\titulo{Teste do Sistema de Listas Automatizadas do PPGT}
\autor{João}{Silva}
\orientador{Prof. Dr. Maria Santos}{UnB}
\datadefesa{15}{12}{2024}
\coordenador{Prof. Dr. Carlos Oliveira}{UnB}

% Membros da banca (obrigatório)
\membrobanca{Prof. Dr. Ana Costa}{UnB}
\membrobanca{Prof. Dr. Pedro Lima}{UFG}

% Palavras-chave para teste
\palavraschave[transportes, listas, LaTeX]{transportes, listas automatizadas, LaTeX}
\keywords[transport, lists, LaTeX]{transport, automated lists, LaTeX}

% Definindo algumas siglas para teste
\sigla{PPGT}{Programa de Pós-Graduação em Transportes}
\sigla{UnB}{Universidade de Brasília}
\sigla{FT}{Faculdade de Tecnologia}
\sigla{ENC}{Departamento de Engenharia Civil e Ambiental}

\begin{document}

% Teste das listas automatizadas
% As listas serão geradas automaticamente se houver conteúdo

\chapter{Introdução}

Este é um teste do sistema de listas automatizadas do PPGT.

% Teste de figura
\begin{figure}[ht]
\centering
\rule{5cm}{3cm} % Placeholder para figura
\caption{Primeira figura de teste para verificar lista automatizada}
\label{fig:teste1}
\end{figure}

% Teste de tabela
\begin{table}[ht]
\centering
\caption{Primeira tabela de teste para verificar lista automatizada}
\label{tab:teste1}
\begin{tabular}{|c|c|}
\hline
Coluna 1 & Coluna 2 \\
\hline
Dados 1 & Dados 2 \\
\hline
\end{tabular}
\end{table}

\chapter{Desenvolvimento}

Mais conteúdo para testar as listas.

% Segunda figura
\begin{figure}[ht]
\centering
\rule{4cm}{2cm} % Placeholder para figura
\caption{Segunda figura de teste}
\label{fig:teste2}
\end{figure}

% Segunda tabela
\begin{table}[ht]
\centering
\caption{Segunda tabela de teste}
\label{tab:teste2}
\begin{tabular}{|l|r|}
\hline
Item & Valor \\
\hline
A & 100 \\
B & 200 \\
\hline
\end{tabular}
\end{table}

Uso de siglas: \gls{PPGT}, \gls{UnB}, \gls{FT}, \gls{ENC}.

\chapter{Conclusão}

As listas foram geradas automaticamente com formatação específica do PPGT.

\end{document}